\section{Many targets for each finite pattern}
\label{sec:targets}

We now fix one of the rounded center lists supplied by
Proposition~\ref{prop:localization}.  Our goal is to bound the probability
that this list covers the window $B_D$ by translates of $K^+$.
The good subset of the window contains many points whose residual vectors
are separated at a small angular scale.  Only a small fraction of pairs
need fail a larger, fixed separation.  The former property controls every
cross diagram; the latter makes most cross diagrams vanish.

\subsection{A deterministic choice of targets}

The following volume estimate controls projective angles, including both
possible orientations of a residual vector.

\begin{lemma}[Cylinder estimate]
\label{lem:cylinder}
Let $n\ge2$, let $D,L>0$, and let $G\subset B_D$ be measurable with
$|G|\ge\omega_nD^n/4$.  Let $Y,Z$ be independent uniform points of $G$,
and let $p,q\in\mathbb R^n$.  For $0<\xi<\pi/2$,
\begin{align}
&\mathbb P\bigl(0<|Y-p|\le L,\ 0<|Z-q|\le L,
                 \ \angle(Y-p,Z-q)\le\xi\bigr)
\notag\\[-2pt]
&\hspace{30mm}\le
 8\frac{\omega_{n-1}}{\omega_n}
       \left(\frac{L\sin\xi}{D}\right)^{n-1}
 \le 8\sqrt n\left(\frac{L\sin\xi}{D}\right)^{n-1}.
\label{eq:cylinder}
\end{align}
\end{lemma}

\begin{proof}
Condition on $Y=y$ with $0<|y-p|\le L$, and write
$u=(y-p)/|y-p|$.  If the event in the statement occurs, the component
of $Z-q$ perpendicular to $u$ has length at most $L\sin\xi$.
Both orientations allowed by the projective angle therefore lie in the
same cylinder around the affine line $q+\mathbb Ru$.
The projection of $B_D$ onto the direction $u$ has length $2D$.
Every perpendicular slice of the cylinder has $(n-1)$-dimensional volume
at most $\omega_{n-1}(L\sin\xi)^{n-1}$.  Thus the allowed points $Z$
have volume at most $2D\omega_{n-1}(L\sin\xi)^{n-1}$.
Divide by $|G|$ and then average over $Y$.

To obtain the last inequality, use the slicing identity
\[
 \frac{\omega_n}{\omega_{n-1}}
 =\int_{-1}^1(1-t^2)^{(n-1)/2}\,dt
 \ge\frac{2\exp(-1/2)}{\sqrt n}.
\]
Indeed, restrict the integral to $|t|\le n^{-1/2}$ and use
$\log(1-1/n)\ge-1/(n-1)$.  In particular,
$\omega_{n-1}/\omega_n\le\sqrt n$.
\end{proof}

For a finite center list $P=(p_1,\ldots,p_N)$ and a point $y$ different
from its entries, write
\[
 \mathcal R_P(y)=(y-p_i:\ |y-p_i|\le b+3\eta)
\]
for its labeled list of relevant residuals.  For $\xi>0$, call a pair
$(y,z)$ \emph{$\xi$-bad} if some residual in $\mathcal R_P(y)$ and
some residual in $\mathcal R_P(z)$ have projective angle at most $\xi$.
This definition is symmetric in $y,z$ and retains repeated center labels.
Recall $\gamma$ from \eqref{eq:rates}, and put
\begin{equation}
 M_n=\left\lfloor\exp\left(\frac{n\log n}{16}\right)\right\rfloor.
\label{eq:target-count}
\end{equation}

\begin{lemma}[Separated targets]
\label{lem:targets}
Fix $H>0$ and $k_0>0$.  For all sufficiently large $n$, let
$P=(p_1,\ldots,p_N)$ be a deterministic list satisfying
\[
 N\le n^{k_0}\exp(\lambda_0n),
\]
and let $G\subset B_D\setminus\{p_1,\ldots,p_N\}$ be measurable with
$|G|\ge\omega_nD^n/4$.
There are $M_n$ distinct points $y_1,\ldots,y_{M_n}\in G$ such that:
\begin{enumerate}[label=\textup{(\roman*)}]
\item no pair with distinct indices is
$H\sqrt{\log n/n}$-bad;
\item at most a fraction $10f_n$ of the unordered pairs with distinct
indices is $4\sqrt\varepsilon$-bad, where
\begin{equation}
 f_n=8\sqrt n\,n^{2k_0}\exp(2\lambda_0n)
       \left(\frac{8\sqrt\varepsilon}{D}\right)^{n-1}
     =\exp\bigl(-\gamma n+O(\log n)\bigr).
\label{eq:bad-pair-rate}
\end{equation}
\end{enumerate}
The constants and the dimension threshold are uniform in $P$ and $G$.
The selected points can be fixed using only $P$ and $G$.
\end{lemma}

\begin{proof}
Sample $M_n$ points independently and uniformly from $G$.
For independent sampled points $Y,Z$, take the union in
Lemma~\ref{lem:cylinder} over all $N^2$ ordered pairs of center labels.
Since $b+3\eta\le2$ for sufficiently large $n$, this gives
\begin{equation}
 \mathbb P\bigl((Y,Z)\text{ is $\xi$-bad}\bigr)
 \le8\sqrt n\,N^2
       \left(\frac{2\sin\xi}{D}\right)^{n-1}.
\label{eq:target-pair}
\end{equation}
All relevant labels are included even though the relevant lists vary
with $Y$ and $Z$.  Repeated centers only repeat events in this union.

Set $t_n=H\sqrt{\log n/n}$.  Multiply the right side of
\eqref{eq:target-pair}, with $\sin t_n\le t_n$, by
$\binom{M_n}{2}$.  The logarithm of the resulting bound is at most
\[
 -\frac38n\log n+\frac12n\log\log n+O_{H,k_0}(n),
\]
which tends to $-\infty$.  Thus the probability that any pair of
sampled points is $t_n$-bad tends to zero, uniformly in the pattern.
For $\xi=4\sqrt\varepsilon$, the inequality
$\sin\xi\le4\sqrt\varepsilon$ instead bounds
\eqref{eq:target-pair} by $f_n$.
The expected fraction of these bad unordered pairs is at most $f_n$.
Markov's inequality shows that this fraction exceeds $10f_n$ with
probability at most $1/10$.
Consequently both conclusions hold with probability at least
$9/10-o(1)>0$.  The sampled points are distinct almost surely because
uniform measure on $G$ has no atoms.  Fix a successful list.
No independence between the two separation properties is needed.
\end{proof}

\subsection{From targets to a small covering probability}

We apply the preceding selection before exposing the Poisson process.
For each selected target, Corollary~\ref{cor:witnesses} supplies sets of
slab normals whose simultaneous occupancy certifies that the target is
uncovered.  Binning their activities will let us combine many certificates
without losing their angular separation.

\begin{proposition}[Failure probability for one pattern]
\label{prop:pattern-tail}
There are constants $c_0,C_{\mathrm{tail}}>0$ and $n_0$ such that, whenever
$n\ge n_0$ and $1\le R\le c_0n\log n$, every retained pattern
$P\in\mathcal F_n$ of Proposition~\ref{prop:localization} satisfies
\begin{equation}
 \mathbb P\left(B_D\subseteq\bigcup_{p\in P}(p+K^+)\right)
 \le \exp\left[-\exp\bigl(\gamma n-C_{\mathrm{tail}}\log n\bigr)\right].
\label{eq:pattern-tail}
\end{equation}
The constants are independent of the pattern, its repetitions, and $R$.
\end{proposition}

The covering event is measurable. Indeed, the finite union of translates
of $K^+$ is closed, so it covers $B_D$ exactly when it contains a fixed
countable dense subset of $B_D$. At each such point, membership is a
finite union of ball conditions and Poisson cap-avoidance events.

\begin{proof}
Fix $P\in\mathcal F_n$ and its deterministic good set $G_P$.
At each $y\in G_P$, Proposition~\ref{prop:localization} gives
\begin{equation}
 r_--2\eta\le |x_i|\le b+3\eta,\qquad
 m_y\le2C_{\mathrm{loc}}R,\qquad
 \sum_{i=1}^{m_y}w_a(|x_i|)\le2C_{\mathrm{loc}}R
\label{eq:target-residual-budget}
\end{equation}
for its relevant residuals $x_i=y-p_i$.
If this list is empty at any point of $G_P$, that point lies in none of
the balls $p+(b+3\eta)B$, so the probability in
\eqref{eq:pattern-tail} is zero.  We may assume all these lists are
nonempty.  For the range of $R$ in the proposition,
$m_y\le n^2$ once $n$ is sufficiently large.

For each relevant residual, define the cutting cap
\[
 C_i(y)=\{u\in S^{n-1}:|u\cdot x_i|>1+3\eta\}.
\]
Lemma~\ref{lem:caps}, applied with radial loss $10\log n/n$
and a fixed perturbation bound, gives
\[
 s_i=-\log\min\{\mu(C_i(y)),1/2\}
       \le C_a w_a(|x_i|).
\]
Consequently
\begin{equation}
 m_y+\sum_{i=1}^{m_y}s_i\le A_0R,
 \qquad A_0=2C_{\mathrm{loc}}(1+C_a).
\label{eq:target-cap-budget}
\end{equation}
All assumptions of Corollary~\ref{cor:witnesses}
therefore hold uniformly in $y$ and $P$.  Fix its constants $A,C_*$ and
$H$ for this value of $A_0$.

The list-length estimate in Proposition~\ref{prop:localization} gives
$N\le n^{k_0}\exp(\lambda_0n)$ for a fixed $k_0$, uniformly in the
present range of $R$.  Apply Lemma~\ref{lem:targets} to $G_P$ with the
just-fixed $H$.  At each of the resulting $M=M_n$ targets, choose the
block partition and measurable witness sets from
Corollary~\ref{cor:witnesses}.  Denote its injection count by $Z_y$,
its activity by $\nu_y$, and the diagram sums by $D_{yz}$.
The corollary gives
\begin{equation}
 \exp(-AR)\le\nu_y\le1,
 \qquad D_{yy}\le\nu_y^2\exp(C_*R).
\label{eq:target-activities}
\end{equation}
Every pair of distinct selected targets satisfies
$D_{yz}\le n^{-3}\nu_y\nu_z$.
Moreover, $D_{yz}=0$ unless the pair is $4\sqrt\varepsilon$-bad.
Our use of ``at most'' in the definition of badness makes both angular
interfaces strict on its complement.

All choices so far depend only on $P$ and its targets.  In particular,
the witness sets are chosen before the Poisson process; their activities
are deterministic intensities.  There is no need to choose witness sets
measurably at every point of $G_P$, since only the finite selected list
is used.  If $Z_y>0$, every relevant residual is cut by a slab of the
process, and the other centers are too far from $y$ to contain it in a
translate of $K^+$.  Thus $Z_y>0$ certifies a hole.

Divide the possible values of $-\log\nu_y$ in $[0,AR]$ into intervals
of length $\log2$.  At most
\[
 Q=1+\left\lceil\frac{AR}{\log2}\right\rceil
\]
bins suffice.  A largest bin has $m\ge M/Q$ targets and, for some
$\nu>0$, all its activities belong to $[\nu,2\nu]$.
There are at most $10f_n\binom M2$ bad unordered pairs in the full
list.  Hence the chosen bin contains at most
\begin{equation}
 10f_nM(M-1)\le10f_nQ^2m^2
\label{eq:bin-bad-pairs}
\end{equation}
bad ordered pairs with distinct indices.
This bound allows the bin to contain every bad pair; no independence
between bin membership and angles is asserted.

Apply Proposition~\ref{prop:poisson} to the witnesses in this bin.
Its numerator is at least $m^2\nu^2$.  By
\eqref{eq:target-activities} and \eqref{eq:bin-bad-pairs}, its denominator
before the factor $2$ is at most
\[
 \sum_{y,z}D_{yz}
 \le4m\nu^2\exp(C_*R)
       +40m^2\nu^2n^{-3}Q^2f_n.
\]
The first term counts diagonal pairs.  Among distinct pairs, only the
coarse bad pairs contribute.  Coverage forces $Z_y=0$ at every target
of the bin, so
\begin{equation}
 \mathbb P\left(B_D\subseteq\bigcup_{p\in P}(p+K^+)\right)
 \le\exp\left[-\frac1{
 8\bigl(\exp(C_*R)/m+10n^{-3}Q^2f_n\bigr)}\right].
\label{eq:bin-tail}
\end{equation}

Now choose
\[
 c_0=\frac1{64\max\{C_*,1\}}.
\]
Uniformly for $1\le R\le c_0n\log n$, we have $\log Q=O(\log n)$ and,
for sufficiently large $n$,
\[
 m\exp(-C_*R)
 \ge\frac{M_n}{Q}\exp(-C_*R)
 \ge\exp\left(\frac{n\log n}{32}\right).
\]
The first term inside the parentheses in \eqref{eq:bin-tail} is
therefore at most $\exp(-n\log n/32)$.  By
\eqref{eq:bad-pair-rate}, the second is
$\exp(-\gamma n+O(\log n))$.
Since $\gamma$ is fixed, their sum is at most
$\exp(-\gamma n+O(\log n))$; absorbing the fixed factor $8$ gives
\eqref{eq:pattern-tail}.

The order of the choices is fixed throughout: the geometric parameters
come first, then the cap and block constants and $H$, then $c_0$, and
finally the dimension threshold.  Every estimate is uniform over the
retained finite class.  In particular, the auxiliary target selection
introduces no dependence on the realization of $\Pi$.
\end{proof}
